\ifdefined\gkver\else\def\gkver{student}\fi
\documentclass[\gkver]{gaokaozhenti}
\begin{document}

\chapter{2015年浙江理综}
%% paperType: 理综物理部分

\begin{enumerate}
\item
%% number: 14
%% paperName: 2015年浙江理综第 14 题 4 分
%% typeId: 1
%% type: 单选题
%% chapter: 电路
%% point: 电容
%% method: 无
%% score: 4
%% degree: 650
%% duplicateId: 0
%% body:
下列说法正确的是\xzanswer{C}

\fourchoices[answer=C]
{电流通过导体的热功率与电流大小成正比}
{力对物体所做的功与力的作用时间成正比}
{电容器所带电荷量与两极板间的电势差成正比}
{弹性限度内，弹簧的劲度系数与弹簧伸长量成正比}

%% answer: C
%% memo:
\memoanswer{
	根据公式 $P=I^2R$ 可知，电阻一定时，电流通过导体的热功率与电流的平方成正比，A 错误；根据公式 $W=Fs\cos\alpha$，力对物体所做的功与力的作用时间无关，B 错误；根据公式 $C=Q/U$，电容器所带电荷量与两极板间的电势差成正比，C 正确；弹簧的劲度系数与弹簧的材料等自身因素有关，D 错误，选项 C 正确。
}

\item
%% number: 15
%% paperName: 2015年浙江理综第 15 题 4 分
%% typeId: 1
%% type: 单选题
%% chapter: 直线运动
%% point: 平均速度和瞬时速度
%% method: 无
%% score: 4
%% degree: 650
%% duplicateId: 0
%% body:
如图所示，气垫导轨上滑块经过光电门时，其上的遮光条将光遮住，电子计时器可自动记录遮光时间 $\Delta t$，测得遮光条的宽度为 $\Delta x$，用 $\frac{\Delta x}{\Delta t}$ 近似代表滑块通过光电门时的瞬时速度，为使 $\frac{\Delta x}{\Delta t}$ 更接近瞬时速度，正确的措施是\xzanswer{A}

\onepicture[h=3.2cm]{figs/15.png}

\fourchoices[answer=A]
{换用宽度更窄的遮光条}
{提高测量遮光条宽度的精确度}
{使滑块的释放点更靠近光电门}
{增大气垫导轨与水平面的夹角}

%% answer: A
%% memo:
\memoanswer{
	平均速度近似等于瞬时速度，即在 $\Delta t\to0$ 的情况下，$\Delta x/\Delta t$ 近似看作瞬时速度，所以要使得 $\Delta x/\Delta t$ 更接近通过光电门的瞬时速度，需要缩短通过时间，即换用宽度更窄的遮光条，选项 A 正确。
}

\item
%% number: 16
%% paperName: 2015年浙江理综第 16 题 4 分
%% typeId: 1
%% type: 单选题
%% chapter: 电场
%% point: 静电
%% method: 无
%% score: 4
%% degree: 650
%% duplicateId: 0
%% body:
如图所示为静电力演示仪，两金属极板分别固定于绝缘支架上，且正对平行放置。工作时两板分别接高压直流电源的正负极，表面镀铝的乒乓球用绝缘细线悬挂在金属极板中间，则\xzanswer{D}

\onepicture[h=3.4cm]{figs/16.png}

\fourchoices[answer=D]
{乒乓球的左侧感应出负电荷}
{乒乓球受到扰动后，会被吸在左极板上}
{乒乓球共受到电场力、重力和库仑力三个力的作用}
{用绝缘棒将乒乓球拨到与右极板接触，放开后乒乓球会在两极板间来回碰撞}

%% answer: D
%% memo:
\memoanswer{
	从图知金属板间电场水平向左，故乒乓球上的电子移到右侧，即乒乓球的右侧感应出负电荷，A 错误；由于静电感应使乒乓球左右两侧感应出的电荷量相等，所以受到的电场力相等，乒乓球受到扰动后，最终仍会静止，不会吸附到左极板上，B 错误；乒乓球受到重力和电场力作用，库仑力即为电场力，C 错误；用绝缘棒将乒乓球拨到与右极板接触使乒乓球带正电，在电场力作用下，运动到左极板与左极板接触，然后乒乓球带负电，接着在电场力作用下，又运动到右极板与右极板接触，乒乓球带正电，在电场力作用下，运动到左极板，如此重复，即乒乓球会在两极板间来回碰撞，选项 D 正确。
}

\item
%% number: 17
%% paperName: 2015年浙江理综第 17 题 4 分
%% typeId: 1
%% type: 单选题
%% chapter: 曲线运动
%% point: 平抛运动
%% method: 无
%% score: 4
%% degree: 450
%% duplicateId: 0
%% body:
如图所示为足球球门，球门宽为 $L$，一个球员在球门中心正前方距离球门 $s$ 处高高跃起，将足球顶入球门的左下方死角（图中 P 点）。球员顶球点的高度为 $h$。足球做平抛运动（足球可看作质点，忽略空气阻力）则\xzanswer{B}

\onepicture[h=3.0cm]{figs/17.png}

\fourchoices[answer=B]
{足球位移大小 $x=\sqrt{\frac{L^{2}}{4}+s^{2}}$}
{足球初速度的大小 $v_{0}=\sqrt{\frac{g}{2h}\left(\frac{L^{2}}{4}+s^{2}\right)}$}
{足球末速度的大小 $v=\sqrt{\frac{g}{2h}\left(\frac{L^{2}}{4}+s^{2}\right)+4gh}$}
{足球初速度的方向与球门线夹角的正切值 $\tan\theta=\frac{L}{2s}$}

%% answer: B
%% memo:
\memoanswer{
	由几何知识得足球的位移 $x=\sqrt{h^2+s^2+\frac{L^{2}}{4}}$，A 错误；足球做平抛运动有 $x'=v_{0}t$，$h=\frac{1}{2}gt^{2}$，解得 $v_{0}=\sqrt{\frac{g}{2h}\left(\frac{L^{2}}{4}+s^{2}\right)}$，B 正确；足球末速度大小为 $v=\sqrt{v_{0}^{2}+(gt)^{2}}=\sqrt{\frac{g}{2h}\left(\frac{L^{2}}{4}+s^{2}\right)+2gh}$，C 错误；足球初速度的方向与球门线夹角的正切值 $\tan\theta=\frac{s}{\frac{L}{2}}=\frac{2s}{L}$，D 错误，选项 B 正确。
}

\item
%% number: 18
%% paperName: 2015年浙江理综第 18 题 6 分
%% typeId: 2
%% type: 多选题
%% chapter: 功和能
%% point: 动能定理
%% method: 无
%% score: 6
%% degree: 650
%% duplicateId: 0
%% body:
我国科学家正在研制航母舰载机使用的电磁弹射器。舰载机总质量为 $3\times10^{4}\Ukg$，设起飞过程中发动机的推力恒为 $1.0\times10^{5}\UN$；弹射器有效作用长度为 $100\Um$，推力恒定。要求舰载机在水平弹射结束时速度大小达到 $80\Ums$。弹射过程中舰载机所受总推力为弹射器和发动机推力之和，假设所受阻力为总推力的 20\%，则\xzanswer{ABD}

\fourchoices[answer=ABD]
{弹射器的推力大小为 $1.1\times10^{6}\UN$}
{弹射器对舰载机所做的功为 $1.1\times10^{8}\UJ$}
{弹射器对舰载机做功的平均功率为 $8.8\times10^{7}\UW$}
{舰载机在弹射过程中的加速度大小为 $32\Umsq$}

%% answer: ABD
%% memo:
\memoanswer{
	设发动机、弹射器的推力分别为 $F_{1}$、$F_{2}$，则阻力 $f=0.2(F_{1}+F_{2})$，由动能定理 $[(F_{1}+F_{2})-f]s=\frac{1}{2}mv^{2}$，而 $F_{1}=1.0\times10^{5}\UN$，故 $F_{2}=1.1\times10^{6}\UN$，A 正确；弹射器对舰载机所做的功为 $W_{2}=F_{2}s=1.1\times10^{8}\UJ$，B 正确；舰载机在弹射过程中的加速度大小为 $a=\frac{(F_{1}+F_{2})-f}{m}=32\Umsq$，D 正确；弹射过程中运动时间 $t=\sqrt{\frac{2s}{a}}=2.5\Us$，弹射器对舰载机做功的平均功率为 $P_{2}=W_{2}/t=4.4\times10^{7}\UW$，C 错误，选项 ABD 正确。
}

\item
%% number: 19
%% paperName: 2015年浙江理综第 19 题 6 分
%% typeId: 2
%% type: 多选题
%% chapter: 曲线运动
%% point: 圆周运动的应用
%% method: 无
%% score: 6
%% degree: 450
%% duplicateId: 0
%% body:
如图所示为赛车场的一个水平“U”形弯道，转弯处为圆心在 O 点的半圆，内外半径分别为 $r$ 和 $2r$。一辆质量为 $m$ 的赛车通过 AB 线经弯道到达 $A'B'$ 线，有如图所示的①②③三条路线，其中路线③是以 $O'$ 为圆心的半圆，$OO'=r$。赛车沿圆弧路线行驶时，路面对轮胎的最大径向静摩擦力为 $F_{\max}$。选择路线，赛车以不打滑的最大速率通过弯道（所选路线内赛车速率不变，发动机功率足够大），则\xzanswer{ACD}

\onepicture[h=4.0cm]{figs/19.png}

\fourchoices[answer=ACD]
{选择路线①，赛车经过的路程最短}
{选择路线②，赛车的速率最小}
{选择路线③，赛车所用时间最短}
{①②③三条路线的圆弧上，赛车的向心加速度大小相等}

%% answer: ACD
%% memo:
\memoanswer{
	由图及几何关系知：路线①的路程为 $s_{1}=2r+\pi r$，路线②的路程为 $s_{2}=2r+2\pi r$，路线③的路程为 $s_{3}=2\pi r$，A 正确；赛车以不打滑的最大速率通过弯道有 $F_{\max}=ma_{n}=m\frac{v^{2}}{R}$，速度 $v=\sqrt{\frac{F_{\max}R}{m}}$，即半径越大，速度越大，选择路线①赛车的速率最小，B 错误、D 正确；根据 $t=s/v$，代入数据解得选择路线③时，赛车所用时间最短，C 正确，选项 ACD 正确。
}

\item
%% number: 20
%% paperName: 2015年浙江理综第 20 题 6 分
%% typeId: 2
%% type: 多选题
%% chapter: 电场
%% point: 带电粒子的平衡
%% method: 无
%% score: 6
%% degree: 450
%% duplicateId: 0
%% body:
如图所示，用两根长度相同的绝缘细线把一个质量为 $0.1\Ukg$ 的小球 A 悬挂到水平板的 MN 两点，A 上带有 $Q=3.0\times10^{-6}\UC$ 的正电荷。两线夹角为 $120^\circ$，两线上的拉力大小分别为 $F_{1}$ 和 $F_{2}$。A 的正下方 $0.3\Um$ 处放有一带等量异种电荷的小球 B，B 与绝缘支架的总质量为 $0.2\Ukg$（重力加速度取 $g=10\Umsq$；静电力常量 $k=9.0\times10^{9}\UNmqCq$，A、B 球可视为点电荷）则\xzanswer{BC}

\onepicture[h=3.0cm]{figs/20.png}

\fourchoices[answer=BC]
{支架对地面的压力大小为 $2.0\UN$}
{两线上的拉力大小 $F_{1}=F_{2}=1.9\UN$}
{将 B 水平右移，使 M、A、B 在同一直线上，此时两线上的拉力大小为 $F_{1}=1.225\UN$，$F_{2}=1.0\UN$}
{将 B 移到无穷远处，两线上的拉力大小 $F_{1}=F_{2}=0.866\UN$}

%% answer: BC
%% memo:
\memoanswer{
	A、B 两电荷的作用力为库仑引力，支架对地面的压力大小为 $F=m_{B}g-k\frac{q_{A}q_{B}}{r^{2}}=1.1\UN$，A 错误；三力平衡夹角恰为 $120^\circ$，则两线上的拉力大小 $F_{1}=F_{2}=m_{A}g+k\frac{q_{A}q_{B}}{r^{2}}=1.9\UN$，B 正确；将 B 水平右移，使 M、A、B 在同一直线上，则两小球距离 $r'=\frac{r}{\sin30^\circ}=0.6\Um$，故 $\left(F_{1}'-k\frac{q_{A}q_{B}}{r'^{2}}\right)\sin60^\circ=F_{2}'\sin60^\circ$，$\left(F_{1}'-k\frac{q_{A}q_{B}}{r'^{2}}\right)\cos60^\circ+F_{2}'\cos60^\circ=m_{A}g$，解得 $F_{1}'=1.225\UN$，$F_{2}'=1.0\UN$，C 正确；将 B 移到无穷远处，两小球的库仑引力为 0，两线上的拉力大小 $F_{1}=F_{2}=m_{A}g=1\UN$，D 错误，选项 BC 正确。
}

\item
%% number: 21
%% paperName: 2015年浙江理综第 21 题 10 分
%% typeId: 4
%% type: 实验题
%% chapter: 直线运动
%% point: 打点计时器公式
%% method: 无
%% score: 10
%% degree: 450
%% duplicateId: 0
%% body:
甲同学准备做“验证机械能守恒定律”实验，乙同学准备做“探究加速度与力、质量的关系”实验。

\onepicture[w=8.5cm]{figs/21a.png}

\begin{enumerate}
	\item
	上图中 A、B、C、D、E 表示部分实验器材，甲同学需在图中选用的器材 \tkanswer[1.6cm]{AB}；乙同学需在图中选用的器材 \tkanswer[1.6cm]{BDE}。（用字母表示）

	\onepicture[w=9.5cm]{figs/21b.png}

	\item
	乙同学在实验室选齐所需器材后，经正确操作获得如图所示的两条纸带①和②。纸带 \tkanswer[1.0cm]{①} 的加速度大（填①或者②），其加速度大小为 \tkanswer[2.6cm]{$(2.5\pm0.2)\Umsq$}。
\end{enumerate}

%% answer:
\jdanswer{
\begin{enumerate}
	\item
	AB；BDE

	\item
	①，$(2.5\pm0.2)\Umsq$
\end{enumerate}
}
%% memo:
\memoanswer{
	（1）“验证机械能守恒定律”实验，需要在竖直面上打出一条重锤下落的纸带，即可验证，故选仪器 AB；“探究加速度与力、质量的关系”实验需要钩码拉动小车打出一条纸带，故选 BDE。

	（2）纸带①中前第 1、2 点与第 2、3 点的位移差为连续相等时间内位移差 $\Delta x_{1}=(32.40-30.70)-(30.70-29.10)=0.1\Ucm$；纸带②中前第 1、2 点与第 2、3 点的位移差为连续相等时间内位移差 $\Delta x_{2}=(31.65-29.00)-(29.00-27.40)=0.05\Ucm$；由匀变速直线运动 $\Delta x=aT^{2}$，可知纸带①加速度大，为 $a=\Delta x/T^{2}=2.5\Umsq$。
}

\item
%% number: 22
%% paperName: 2015年浙江理综第 22 题 10 分
%% typeId: 4
%% type: 实验题
%% chapter: 电路
%% point: 伏安法测电阻
%% method: 无
%% score: 10
%% degree: 450
%% duplicateId: 0
%% body:
图 1 是小红同学在做“描绘小灯泡的伏安特性曲线”实验的实物连接图。

\twopicture[numA={1}, numB={2}, hA=5.0cm, hB=5.0cm]
{figs/22a.png}
{figs/22b.png}

\begin{enumerate}
	\item
	根据图 1 画出实验电路图。

	\item
	调节滑动变阻器得到了两组电流表与电压表的示数如图 2 中的①②③④所示，电流表量程为 $0.6\UA$，电压表量程为 $3\UV$。所示读数为：①\tkanswer[1.4cm]{0.10}$\UA$、②\tkanswer[1.4cm]{0.24}$\UA$、③\tkanswer[1.4cm]{2.00}$\UV$、④\tkanswer[1.4cm]{0.27}$\UV$。两组数据得到的电阻分别为 \tkanswer[2.4cm]{$(8.3\pm0.1)\UO$} 和 \tkanswer[2.4cm]{$(2.7\pm0.1)\UO$}。
\end{enumerate}

%% answer:
\jdanswer{
\begin{enumerate}
	\item
	如图所示。

	\item
	①$0.10\UA$、②$0.24\UA$、③$2.00\UV$、④$0.27\UV$；$(8.3\pm0.1)\UO$ 和 $(2.7\pm0.1)\UO$。
\end{enumerate}
}
%% memo:
\memoanswer{
	（1）从实物图中可知电压表测量灯泡两端电压，电流表采用外接法，滑动变阻器采用分压接法，故电路图如图所示。

	\onepicture[h=3.2cm, align=left]{figs/22ans.png}

	（2）电流表的量程为 $0.6\UA$，最小分度值为 $0.02\UA$，①的读数为 $0.10\UA$，②的读数为 $0.24\UA$；电压表量程为 $3\UV$，最小分度值为 $0.1\UV$，③的读数为 $2.00\UV$，④的读数为 $0.27\UV$。由 $R=U/I$ 知：①④为同组数据得电阻为 $2.7\UO$，②③为另组数据得电阻为 $8.3\UO$。
}

\item
%% number: 23
%% paperName: 2015年浙江理综第 23 题 18 分
%% typeId: 6
%% type: 计算题
%% chapter: 功和能
%% point: 动能定理
%% method: 无
%% score: 18
%% degree: 450
%% duplicateId: 0
%% body:
如图所示，用一块长 $L_{1}=1.0\Um$ 的木板在墙和桌面间架设斜面，桌面高 $H=0.8\Um$，长 $L_{2}=1.5\Um$。斜面与水平桌面的倾角 $\theta$ 可在 $0\sim60^\circ$ 间调节后固定。将质量 $m=0.2\Ukg$ 的小物块从斜面顶端静止释放，物块与斜面间的动摩擦因数 $\mu_{1}=0.05$，物块与桌面间的动摩擦因数 $\mu_{2}$，忽略物块在斜面与桌面交接处的能量损失。（重力加速度取 $g=10\Umsq$；最大静摩擦力等于滑动摩擦力）

\begin{enumerate}
	\item
	求 $\theta$ 角增大到多少时，物块能从斜面开始下滑；（用正切值表示）

	\item
	当 $\theta$ 增大到 $37^\circ$ 时，物块恰能停在桌面边缘，求物块与桌面间的动摩擦因数 $\mu_{2}$（已知 $\sin37^\circ=0.6$，$\cos37^\circ=0.8$）

	\item
	继续增大 $\theta$ 角，发现 $\theta=53^\circ$ 时物块落地点与墙面的距离最大，求此最大距离 $x_{m}$。
\end{enumerate}

\onepicture[h=3.6cm, align=right]{figs/23.png}

%% answer:
\jdanswer{
\begin{enumerate}
	\item
	$\tan\theta\geq0.05$

	\item
	$\mu_{2}=0.8$

	\item
	$x_{m}=1.9\Um$
\end{enumerate}
}
%% memo:
\memoanswer{
	（1）为使小物块下滑 $mg\sin\theta\geq\mu_{1}mg\cos\theta$，$\theta$ 满足的条件 $\tan\theta\geq0.05$。

	（2）克服摩擦力做功 $W_{f}=\mu_{1}mgL_{1}\cos\theta+\mu_{2}mg(L_{2}-L_{1}\cos\theta)$，由动能定理得 $mgL_{1}\sin\theta-W_{f}=0$，代入数据得 $\mu_{2}=0.8$。

	（3）由动能定理得 $mgL_{1}\sin\theta-W_{f}=\frac{1}{2}mv^{2}$，代入数据得 $v=1\Ums$，$H=\frac{1}{2}gt^{2}$，$t=0.4\Us$，$x_{1}=vt=0.4\Um$，$x_{m}=x_{1}+L_{2}=1.9\Um$。
}

\item
%% number: 24
%% paperName: 2015年浙江理综第 24 题 20 分
%% typeId: 6
%% type: 计算题
%% chapter: 磁场
%% point: 安培力
%% method: 无
%% score: 20
%% degree: 450
%% duplicateId: 0
%% body:
小明同学设计了一个“电磁天平”，如图 1 所示，等臂天平的左臂为挂盘，右臂挂有矩形线圈，两臂平衡。线圈的水平边长 $L=0.1\Um$，竖直边长 $H=0.3\Um$，匝数为 $N_{1}$。线圈的下边处于匀强磁场内，磁感应强度 $B_{0}=1\UT$，方向垂直线圈平面向里。线圈中通有可在 $0\sim2.0\UA$ 范围内调节的电流 $I$。挂盘放上待测物体后，调节线圈中电流使得天平平衡，测出电流即可测得物体的质量。（重力加速度取 $g=10\Umsq$）

\begin{enumerate}
	\item
	为使电磁天平的量程达到 $0.5\Ukg$，线圈的匝数 $N_{1}$ 至少为多少

	\item
	进一步探究电磁感应现象，另选 $N_{2}=100$ 匝、形状相同的线圈，总电阻 $R=10\UO$，不接外电流，两臂平衡，如图 2 所示，保持 $B_{0}$ 不变，在线圈上部另加垂直纸面向外的匀强磁场，且磁感应强度 $B$ 随时间均匀变大，磁场区域宽度 $d=0.1\Um$。当挂盘中放质量为 $0.01\Ukg$ 的物体时，天平平衡，求此时磁感应强度的变化率 $\frac{\Delta B}{\Delta t}$。
\end{enumerate}

\twopicture[numA={1}, numB={2}, hA=5.0cm, hB=5.0cm, align=right]
{figs/24a.png}
{figs/24b.png}

%% answer:
\jdanswer{
\begin{enumerate}
	\item
	$N_{1}=25$ 匝

	\item
	$0.1\UTs$
\end{enumerate}
}
%% memo:
\memoanswer{
	（1）线圈受到安培力 $F=N_{1}B_{0}IL$，天平平衡 $mg=N_{1}B_{0}IL$，代入数据得 $N_{1}=25$ 匝。

	（2）由电磁感应定律得 $E=N_{2}\frac{\Delta\Phi}{\Delta t}=N_{2}\frac{\Delta B}{\Delta t}Ld$，由欧姆定律得 $I'=\frac{E}{R}$，线圈受到安培力 $F'=N_{2}B_{0}I'L$，天平平衡 $m'g=N_{2}^{2}B_{0}\frac{\Delta B}{\Delta t}\cdot\frac{dL^{2}}{R}$，代入数据可得 $\frac{\Delta B}{\Delta t}=0.1\UTs$。
}

\item
%% number: 25
%% paperName: 2015年浙江理综第 25 题 22 分
%% typeId: 6
%% type: 计算题
%% chapter: 磁场
%% point: 带电粒子在磁场中的运动
%% method: 无
%% score: 22
%% degree: 450
%% duplicateId: 0
%% body:
使用回旋加速器的实验需要把离子束从加速器中引出，离子束引出的方法有磁屏蔽通道法和静电偏转法等。质量为 $m$、速度为 $v$ 的离子在回旋加速器内旋转，旋转轨道是半径为 $r$ 的圆，圆心在 O 点，轨道在垂直纸面向外的匀强磁场中，磁感应强度为 $B$。

为引出离子束，使用磁屏蔽通道法设计引出器。引出器原理如图所示，一对圆弧形金属板组成弧形引出通道，通道的圆心位于 $O'$ 点（$O'$ 点图中未画出）。引出离子时，令引出通道内磁场的磁感应强度降低，从而使离子从 P 点进入通道，沿通道中心线从 Q 点射出。已知 OQ 长度为 $L$，OQ 与 OP 的夹角为 $\theta$，

\begin{enumerate}
	\item
	求离子的电荷量 $q$ 并判断其正负；

	\item
	离子从 P 点进入，Q 点射出，通道内匀强磁场的磁感应强度应降为 $B'$，求 $B'$；

	\item
	换用静电偏转法引出离子束，维持通道内的原有磁感应强度 $B$ 不变，在内外金属板间加直流电压，两板间产生径向电场，忽略边缘效应。为使离子仍从 P 点进入，Q 点射出，求通道内引出轨迹处电场强度 $E$ 的方向和大小。
\end{enumerate}

\onepicture[h=3.8cm, align=right]{figs/25.png}

%% answer:
\jdanswer{
\begin{enumerate}
	\item
	$q=\frac{mv}{Br}$，正电荷

	\item
	$\frac{mv(2r-2L\cos\theta)}{q(r^{2}+L^{2}-2rL\cos\theta)}$

	\item
	$E=Bv-\frac{mv^{2}(2r-2L\cos\theta)}{q(r^{2}+L^{2}-2rL\cos\theta)}$
\end{enumerate}
}
%% memo:
\memoanswer{
	（1）离子做圆周运动 $Bqv=\frac{mv^{2}}{r}$，$q=\frac{mv}{Br}$，正电荷。

	（2）如图所示，$O'Q=R$，$OQ=L$，$O'O=R-r$，引出轨迹为圆弧 $B'qv=\frac{mv^{2}}{R}$，$R=\frac{mv}{qB'}$。

	根据几何关系得 $R=\frac{r^{2}+L^{2}-2rL\cos\theta}{2r-2L\cos\theta}$，$B'=\frac{mv}{qR}=\frac{mv(2r-2L\cos\theta)}{q(r^{2}+L^{2}-2rL\cos\theta)}$。

	\onepicture[h=4.0cm, align=left]{figs/25sol.png}

	（3）电场强度方向沿径向向外，引出轨迹为圆弧 $Bqv-Eq=\frac{mv^{2}}{R}$，$E=Bv-\frac{mv^{2}(2r-2L\cos\theta)}{q(r^{2}+L^{2}-2rL\cos\theta)}$。
}

\end{enumerate}

\end{document}
